%! Change in Linear and Exponential Functions — Unit 4, Lesson 4.2 — Warm-Up
\documentclass[10pt]{article}
\usepackage{precalculus-article}
\usepackage{precalculus-boxes}
\newcommand{\TallMath}[1]{$\displaystyle #1\rule[-1.4em]{0pt}{3.2em}$}

\begin{document}

\pageheader{Unit 4, Lesson 4.2}{Warm-Up}
\namedateperiod

\vspace{0.14in}

\begin{enumerate}[itemsep=18pt]

\item The sequence $5,\ 8,\ 11,\ 14,\ \ldots$ is arithmetic.

      \begin{enumerate}[label=(\alph*), itemsep=8pt]
        \item Find the common difference $d$. \quad $d = $ \blank{3cm}

        \item Write the explicit formula $a_n$ in terms of $n$ (starting at $n = 0$). \\[4pt]
              $a_n = $ \blank{6cm}
      \end{enumerate}

\item The sequence $4,\ 12,\ 36,\ 108,\ \ldots$ is geometric.

      \begin{enumerate}[label=(\alph*), itemsep=8pt]
        \item Find the common ratio $r$. \quad $r = $ \blank{3cm}

        \item Write the explicit formula $g_n$ in terms of $n$ (starting at $n = 0$). \\[4pt]
              $g_n = $ \blank{6cm}
      \end{enumerate}

\item A sequence has $a_3 = 20$ and $a_7 = 36$.

      \begin{enumerate}[label=(\alph*), itemsep=8pt]
        \item Is the sequence arithmetic or geometric? Circle one and show your check.

              \medskip
              \textbf{Arithmetic} \quad / \quad \textbf{Geometric}

              \vspace{0.08in}
              Check: \blank{8cm}

        \item Find $a_0$. \quad $a_0 = $ \blank{3cm}
      \end{enumerate}

\end{enumerate}

\end{document}
